% GENERATED FILE. Edit canonical lessons, then run npm run generate:books.
% canonical-source-hash: fnv1a64:67d86ae283d6c5da
% canonical-lessons: GU-C517-vayar, GU-C517-plag, GU-C517-svic, GU-C517-balb, GU-C517-khosvu

\chapter{Wires and Switches}
\label{ch:gu-wires-and-switches}
\hlchaptermodality{517}

\begin{chapteropening}
\textbf{By the end of this chapter:} \emph{I can talk about wires, plugs, switches and bulbs.}

Complete the last lesson of chapter 517: I can talk about wires, plugs, switches and bulbs.
\end{chapteropening}

\section[\texorpdfstring{\gu{વાયર} (vāyar)}{વાયર (vāyar)}]{\gu{વાયર} (vāyar) — an electric wire}

\label{lesson:GU-C517-vayar}



\begin{quote}
Before the new one: say the Gujarati for to colour, then the Gujarati for to plaster.
\end{quote}

\subsection*{You'll want to know: \gu{વાયર}}
\textbf{\gu{વાયર}} — \emph{vāyar} — \textquotedblleft{}an electric wire\textquotedblright{}.

\begin{quote}
\gu{પંખાનો} \gu{વાયર} \gu{જૂનો} \gu{થઈ} \gu{ગયો} \gu{છે}.
\end{quote}

\begin{quote}
\emph{pankhāno vāyar jūno thaī gayo che.}
\end{quote}

\begin{quote}
The fan's wire has got old.
\end{quote}

A masculine noun.

\begin{cousinweb}
\textbf{\gu{વાયર}} (\emph{vāyar}) is borrowed from English \emph{wire}.
\end{cousinweb}

\begin{grammarlens}[title={What you've built}]
One more word for describing what you lived and felt.
\end{grammarlens}

\subsection*{Your turn}
\begin{itemize}
\raggedright
  \item \emph{Say it:} \emph{vāyar}
  \item \emph{Say it:} \emph{vāyar}, once more
  \item \emph{Say it:} \emph{pankhāno vāyar jūno thaī gayo che.} — The fan's wire has got old.
  \item \emph{Recall:} say the Gujarati for to colour, then the Gujarati for to plaster, then say \emph{vāyar} again
\end{itemize}

\begin{tcolorbox}[breakable,colback=teal!4,colframe=teal!35!black,title={Before you move on}]
What does \gu{વાયર} mean? (\textquotedblleft{}An electric wire\textquotedblright{}.) Say it once more.
\end{tcolorbox}
\section[\texorpdfstring{\gu{પ્લગ} (plag)}{પ્લગ (plag)}]{\gu{પ્લગ} (plag) — a plug}

\label{lesson:GU-C517-plag}



\begin{quote}
Before the new one: say the Gujarati for to plaster, then the Gujarati for an electric wire.
\end{quote}

\subsection*{You'll want to know: \gu{પ્લગ}}
\textbf{\gu{પ્લગ}} — \emph{plag} — \textquotedblleft{}a plug\textquotedblright{}.

\begin{quote}
\gu{ફોનનો} \gu{પ્લગ} \gu{ક્યાં} \gu{છે}?
\end{quote}

\begin{quote}
\emph{phonno plag kyā̃ che?}
\end{quote}

\begin{quote}
Where is the phone's plug?
\end{quote}

A masculine noun.

\begin{cousinweb}
\textbf{\gu{પ્લગ}} (\emph{plag}) is borrowed from English \emph{plug}.
\end{cousinweb}

\begin{grammarlens}[title={What you've built}]
One more word for describing what you lived and felt.
\end{grammarlens}

\subsection*{Your turn}
\begin{itemize}
\raggedright
  \item \emph{Say it:} \emph{plag}
  \item \emph{Say it:} \emph{plag}, once more
  \item \emph{Say it:} \emph{phonno plag kyā̃ che?} — Where is the phone's plug?
  \item \emph{Recall:} say the Gujarati for to plaster, then the Gujarati for an electric wire, then say \emph{plag} again
\end{itemize}

\begin{tcolorbox}[breakable,colback=teal!4,colframe=teal!35!black,title={Before you move on}]
What does \gu{પ્લગ} mean? (\textquotedblleft{}A plug\textquotedblright{}.) Say it once more.
\end{tcolorbox}
\section[\texorpdfstring{\gu{સ્વિચ} (svic)}{સ્વિચ (svic)}]{\gu{સ્વિચ} (svic) — a switch}

\label{lesson:GU-C517-svic}



\begin{quote}
Before the new one: say the Gujarati for an electric wire, then the Gujarati for a plug.
\end{quote}

\subsection*{You'll want to know: \gu{સ્વિચ}}
\textbf{\gu{સ્વિચ}} — \emph{svic} — \textquotedblleft{}a switch\textquotedblright{}.

\begin{quote}
\gu{બારણા} \gu{પાસે} \gu{બત્તીની} \gu{સ્વિચ} \gu{છે}.
\end{quote}

\begin{quote}
\emph{bārṇā pāse battīnī svic che.}
\end{quote}

\begin{quote}
The light switch is by the door.
\end{quote}

A feminine noun.

\begin{cousinweb}
\textbf{\gu{સ્વિચ}} (\emph{svic}) is borrowed from English \emph{switch}.
\end{cousinweb}

\begin{grammarlens}[title={What you've built}]
One more word for describing what you lived and felt.
\end{grammarlens}

\subsection*{Your turn}
\begin{itemize}
\raggedright
  \item \emph{Say it:} \emph{svic}
  \item \emph{Say it:} \emph{svic}, once more
  \item \emph{Say it:} \emph{bārṇā pāse battīnī svic che.} — The light switch is by the door.
  \item \emph{Recall:} say the Gujarati for an electric wire, then the Gujarati for a plug, then say \emph{svic} again
\end{itemize}

\begin{tcolorbox}[breakable,colback=teal!4,colframe=teal!35!black,title={Before you move on}]
What does \gu{સ્વિચ} mean? (\textquotedblleft{}A switch\textquotedblright{}.) Say it once more.
\end{tcolorbox}
\section[\texorpdfstring{\gu{બલ્બ} (balb)}{બલ્બ (balb)}]{\gu{બલ્બ} (balb) — a light bulb}

\label{lesson:GU-C517-balb}



\begin{quote}
Before the new one: say the Gujarati for a plug, then the Gujarati for a switch.
\end{quote}

\subsection*{You'll want to know: \gu{બલ્બ}}
\textbf{\gu{બલ્બ}} — \emph{balb} — \textquotedblleft{}a light bulb\textquotedblright{}.

\begin{quote}
\gu{ઓરડાનો} \gu{બલ્બ} \gu{ઊડી} \gu{ગયો}.
\end{quote}

\begin{quote}
\emph{orḍāno balb ūḍī gayo.}
\end{quote}

\begin{quote}
The bulb in the room has blown.
\end{quote}

A masculine noun.

\begin{cousinweb}
\textbf{\gu{બલ્બ}} (\emph{balb}) is borrowed from English \emph{bulb}. A bulb that blows \textquotedblleft{}flies off\textquotedblright{}: \textbf{\gu{બલ્બ} \gu{ઊડી} \gu{ગયો}}.
\end{cousinweb}

\begin{grammarlens}[title={What you've built}]
One more word for describing what you lived and felt.
\end{grammarlens}

\subsection*{Your turn}
\begin{itemize}
\raggedright
  \item \emph{Say it:} \emph{balb}
  \item \emph{Say it:} \emph{balb}, once more
  \item \emph{Say it:} \emph{orḍāno balb ūḍī gayo.} — The bulb in the room has blown.
  \item \emph{Recall:} say the Gujarati for a plug, then the Gujarati for a switch, then say \emph{balb} again
\end{itemize}

\begin{tcolorbox}[breakable,colback=teal!4,colframe=teal!35!black,title={Before you move on}]
What does \gu{બલ્બ} mean? (\textquotedblleft{}A light bulb\textquotedblright{}.) Say it once more.
\end{tcolorbox}
\section[\texorpdfstring{\gu{ખોસવું} (khosvũ)}{ખોસવું (khosvũ)}]{\gu{ખોસવું} (khosvũ) — to push in (firmly)}

\label{lesson:GU-C517-khosvu}



\begin{quote}
Before the new one: say the Gujarati for a switch, then the Gujarati for a light bulb.
\end{quote}

\subsection*{You'll want to know: \gu{ખોસવું}}
\textbf{\gu{ખોસવું}} — \emph{khosvũ} — \textquotedblleft{}to push in (firmly)\textquotedblright{}.

\begin{quote}
\gu{અરુણે} \gu{પ્લગ} \gu{દીવાલમાં} \gu{ખોસ્યો}.
\end{quote}

\begin{quote}
\emph{Aruṇe plag dīvālmā̃ khosyo.}
\end{quote}

\begin{quote}
Arun pushed the plug into the wall.
\end{quote}

Past, with \gu{તેણે} (he or she …): \gu{તેણે} \gu{પ્લગ} \gu{ખોસ્યો}, \emph{teṇe plag khosyo}.

\begin{cousinweb}
\textbf{\gu{ખોસવું}} (\emph{khosvũ}) comes from \emph{*skōṣati}, a form CDIAL reconstructs, and CDIAL lists the Gujarati verb as \textquotedblleft{}to force in\textquotedblright{} (CDIAL 13661).
\end{cousinweb}

\begin{grammarlens}[title={What you've built}]
One more word for describing what you lived and felt.
\end{grammarlens}

\subsection*{Your turn}
\begin{itemize}
\raggedright
  \item \emph{Say it:} \emph{khosvũ}
  \item \emph{Say it:} \emph{khosvũ}, once more
  \item \emph{Say it:} \emph{Aruṇe plag dīvālmā̃ khosyo.} — Arun pushed the plug into the wall.
  \item \emph{Recall:} say the Gujarati for a switch, then the Gujarati for a light bulb, then say \emph{khosvũ} again
\end{itemize}

\begin{tcolorbox}[breakable,colback=teal!4,colframe=teal!35!black,title={Before you move on}]
What does \gu{ખોસવું} mean? (\textquotedblleft{}To push in (firmly)\textquotedblright{}.) Say it once more.
\end{tcolorbox}
